% 仙后座（综述）
% license CCBYSA3
% type Wiki

本文根据 CC-BY-SA 协议转载翻译自维基百科\href{https://en.wikipedia.org/wiki/Cassiopeia_(constellation)}{相关文章}。

仙后座（Cassiopeia，ⓘ）是位于北天的一個星座與星群，其名稱來自希臘神話中虛榮的王后卡西奧佩婭（Cassiopeia），她是安德洛墨達（Andromeda）的母親，曾誇耀自己擁有無與倫比的美貌。仙后座是公元 2 世紀希臘天文學家托勒密列出的 48 個星座之一，至今仍是現代 88 個星座之一。由於其由五顆亮星組成的獨特「W」形狀，仙后座極易辨認。

仙后座位於北天，在北緯 34° 以上的地區全年可見；在（亞）熱帶地區，9 月至 11 月初觀測條件最佳；而在南半球低緯度、低於南緯 25° 的熱帶地區，則可於特定季節在北方低空看到。

視星等為 2.2 的仙后座 α 星，即舍達爾（Schedar），是仙后座中最亮的恆星。該星座擁有一些已知最為明亮的恆星，包括黃色特超巨星仙后座 ρ 星（Rho Cassiopeiae）與 V509 Cassiopeiae，以及白色特超巨星仙后座 6 星。1572 年，第谷·布拉赫（Tycho Brahe）所觀測到的超新星就在仙后座中明亮爆發$^{[5]}$。仙后座 A 是一個超新星殘骸，在頻率高於 1 GHz 時，是全天最亮的銀河系外射電源。已有 14 個恆星系統被發現擁有系外行星，其中之一——HD 219134——被認為可能擁有六顆行星。一段富集的銀河帶穿過仙后座，其中包含多個疏散星團、年輕而明亮的銀河盤恆星以及星雲。IC 10 是一個不規則星系，是目前已知距離最近的星暴星系，也是本星系群中唯一的此類星系。
\subsection{神話}
\begin{figure}[ht]
\centering
\includegraphics[width=6cm]{./figures/99da126f8a82dd17.png}
\caption{《乌拉尼娅之镜》中描绘的坐在宝座上的仙后座（Cassiopeia）} \label{fig_Cape_1}
\end{figure}
该星座以埃塞俄比亚王后卡西奥佩娅（Cassiopeia）命名。卡西奥佩娅是埃塞俄比亚国王刻甫斯（Cepheus）的妻子$^{[5]}$，也是公主安德洛墨达（Andromeda）的母亲。刻甫斯与卡西奥佩娅以及安德洛墨达一同被安置在星空中，相互毗邻。她之所以被放置在天空中，作为惩罚，是因为她夸口说自己的女儿安德洛墨达比涅瑞伊得斯（Nereids，海中女神）更加美丽，或者另一种说法是，她自称比海中仙女本身更加美丽，从而激怒了波塞冬$^{[6]}$。她被迫坐在宝座上绕北天极旋转，一半时间必须紧紧抓住宝座以免跌落；而波塞冬则裁定安德洛墨达应被绑在岩石上，作为海怪刻托斯（Cetus）的祭品。随后，安德洛墨达被英雄珀尔修斯所拯救，并最终与其成婚$^{[7][8]}$。

在历史上，卡西奥佩娅星座的形象有多种不同的描绘方式。在波斯，她被苏菲（al-Sufi）描绘为一位女王，右手持有一根带新月的权杖，头戴王冠，身旁还有一只双峰骆驼；在法国，她则被描绘为坐在大理石宝座上，左手持棕榈叶，右手提着自己的长袍。这一形象出自奥古斯丁·罗耶（Augustin Royer）于1679年出版的星图集$^{[7]}$。

在中国天文学中，构成仙后座的恒星分布于三个区域之中：紫微垣（紫微垣，Zǐ Wēi Yuán）、北方玄武（北方玄武，Běi Fāng Xuán Wǔ）以及西方白虎（西方白虎，Xī Fāng Bái Hǔ）。

中国天文学家在现代所称的仙后座区域内辨认出多种星象组合。κ、η 与 μ 仙后座星组成了名为“王桥”的星官；当它们与 α 与 β 仙后座星共同观测时，则构成了名为“王良”的大车。御者手中的鞭子由 γ 仙后座星表示，该星有时被称为“次”，即汉语中“鞭”的意思$^{[7]}$。

在印度神话中，仙后座（Cassiopeia）与神话人物沙尔米什塔（Sharmishtha）相关联——她是伟大的阿修罗（Daitya）国王弗里什帕尔瓦（Vrishparva）的女儿，也是德瓦雅尼（Devayani，对应安德洛墨达）的朋友。

在威尔士神话中，Llys Dôn（字面意思为“多恩的宫廷”）是该星座的传统威尔士名称。多恩的至少三位子女同样具有天文学上的对应：Caer Gwydion（“圭迪恩的堡垒”）是银河的传统威尔士名称，而 Caer Arianrhod（“阿丽安罗德的堡垒”）则对应北冕座（Corona Borealis）$^{[9]}$。

在17世纪，仙后座中的恒星曾被用来描绘多位《圣经》人物，其中包括所罗门之母拔示巴（Bathsheba）、《旧约》中的女先知底波拉（Deborah），以及耶稣的追随者抹大拉的玛利亚（Mary Magdalene）$^{[7]}$。

在一些阿拉伯星图集中，仙后座的恒星还被描绘为一个名为“染色之手”（Tinted Hand）的形象。据不同说法，这只手要么象征一只被指甲花染成红色的女子之手，要么象征穆罕默德之女法蒂玛（Fatima）沾满鲜血的手。该“手”由 α、β、γ、δ、ε 与 η 仙后座星组成；其“手臂”则由 α、γ、δ、ε、η 与 ν 英仙座星组成$^{[7]}$。

另一个包含仙后座恒星的阿拉伯星座是“骆驼”。其头部由仙女座的 λ、κ、ι 与 φ 星构成；驼峰是 β 仙后座星；身体由其余仙后座恒星组成；而双腿则由英仙座与仙女座中的恒星构成$^{[7]}$。

在其他文化中，人们还从这一星群中看到了手或麋鹿鹿角的形象$^{[10]}$。例如在萨米人（Sámi）的传统中，仙后座的 “W” 形状被视为麋鹿的鹿角；西伯利亚的楚科奇人（Chukchi）同样将其中五颗主星视作五只驯鹿雄鹿$^{[7]}$。

马绍尔群岛的居民将仙后座视为一幅巨大鼠海豚星座的一部分：仙后座的主星构成其尾部，仙女座与三角座（Triangulum）构成其身体，而白羊座则构成其头部$^{[7]}$。在夏威夷，α、β 与 γ 仙后座星各自有名称：α 仙后座星被称为 Poloahilani，β 仙后座星被称为 Polula，γ 仙后座星被称为 Mulehu。普卡普卡（Pukapuka）群岛的人们则将仙后座视为一个独立的星座，称为 Na Taki-tolu-a-Mataliki $^{[11]}$。
\subsection{特征}
\begin{figure}[ht]
\centering
\includegraphics[width=6cm]{./figures/8d913297d824a2aa.png}
\caption{夜空中的仙后座} \label{fig_Cape_2}
\end{figure}
仙后座曾发生过一次超新星事件，即仙后座A（Cassiopeia A），SN~1572。

仙后座覆盖面积为 $598.4$ 平方度，占全天面积的 $1.451\%$，在 $88$ 个星座中按面积排名第 $25$ 位$^{[12]}$。它北侧和西侧与仙王座（Cepheus）相邻，南侧和西侧与仙女座（Andromeda）相邻，东南侧与英仙座（Perseus）相邻，东侧与鹿豹座（Camelopardalis）相邻，并在西侧与蝎虎座（Lacerta）共享一小段边界。

该星座的三字母官方缩写由国际天文学联合会（International Astronomical Union）于 1922 年采用，为 ``Cas''$^{[13]}$。官方星座边界由比利时天文学家欧仁·德尔波特（Eugène Delporte）于 $1930$ 年制定$^{[b]}$，以一个由 $30$ 条边组成的多边形加以界定。在赤道坐标系中，这些边界的赤经范围为 $00^{\mathrm h}\,27^{\mathrm m}\,03^{\mathrm s}$ 至 $23^{\mathrm h}\,41^{\mathrm m}\,06^{\mathrm s}$，赤纬范围为 $+77.69^\circ$ 至 $+46.68^\circ$,$^{[3]}$。由于其位于北天球，该星座对南纬 $12^\circ$ 以北的观测者而言均可见$^{[12][c]}$。在高纬北方天空中，它属于拱极星座（即在夜空中永不落下），在不列颠群岛、加拿大以及美国北部地区均可全年观测到$^{[15]}$。
\subsubsection{特征}
\textbf{恒星}
\begin{figure}[ht]
\centering
\includegraphics[width=6cm]{./figures/88f670490ee393b8.png}
\caption{从北半球观测时，肉眼所见的仙后座星座} \label{fig_Cape_3}
\end{figure}
德国制图学家约翰·拜耳（Johann Bayer）使用希腊字母 Alpha 至 Omega，随后又使用 A 和 B，来标记该星座中最显著的 26 颗恒星。后来发现 Upsilon 实际上由两颗恒星组成，并由约翰·弗兰斯蒂德（John Flamsteed）标记为 Upsilon$^{[1]}$ 和 Upsilon$^{[2]}$。B Cassiopeiae 实际上就是著名的超新星——第谷超新星（Tycho's Supernova）$^{[16]}$。在该星座边界之内，共有 157 颗视星等小于或等于 6.5 的恒星$^{[d][12]}$。

“W”星群  
仙后座中最亮的五颗恒星——Alpha、Beta、Gamma、Delta 与 Epsilon Cassiopeiae——构成了具有代表性的 W 形星群$^{[15]}$。这五颗恒星均可用肉眼清晰观测，其中三颗具有明显的变光性，第四颗则被怀疑为低振幅变星。在北半球的春季与夏季夜晚，当该星群位于北极星下方时，其形状呈现为 W；而在北半球冬季，或从南半球纬度观测时，该星群位于北极星“上方”（即更接近天顶），其 W 形看起来则是倒置的。

Alpha Cassiopeiae，传统名称为 Schedar（源自阿拉伯语 Al Sadr，意为“胸部”），常被误认为是一个由四颗恒星组成的系统，但实际上它是一颗主星，伴随三颗在物理上距离遥远的视伴星。主星占据主导地位，是一颗橙色巨星，视星等为 2.2，距离地球$228\pm2$光年$^{[18]}$。其光度约为太阳的 771 倍，在其约 1 至 2 亿年的生命周期中，核心氢燃料耗尽后发生膨胀并冷却；在此前的大部分时间里，它是一颗蓝白色 B 型主序星$^{[19]}$。一颗视星等 8.9 的黄色矮星伴星（B）与其相距甚远，而另外两颗伴星（C 与 D）距离更近，视星等分别为 13 和 14$^{[20]}$。

Beta Cassiopeiae，又称 Caph（意为“手”），是一颗白色恒星，视星等为 2.3，距离地球 $54.7\pm0.3$ 光年$^{[18]}$。其年龄约为 12 亿年，核心氢已耗尽，正开始偏离主序阶段并发生膨胀和冷却。其质量约为太阳的 1.9 倍，光度约为太阳的 21.3 倍。Caph 的自转速度约为其临界速度的 92\%，完成一次自转仅需 1.12 天，这使得恒星呈现为扁球体形状，其赤道半径比极半径大约 24\%$^{[21]}$。它是一颗盾牌座$\delta$ 型变星，振幅较小，周期约为 2.5 小时$^{[22]}$。

Gamma Cassiopeiae，现称为 Tiansi，是 Gamma Cassiopeiae 型变星的原型，这类变星由于高速自转，会将物质抛射形成一个可变的盘状结构。Gamma Cassiopeiae 的最暗视星等为 3.0，最亮可达 1.6，但通常接近 2.2，其亮度变化具有不可预测的减弱与增强。它是一个分光双星系统，轨道周期为 203.59 天，其伴星的计算质量约与太阳相当。然而至今尚未获得该伴星的直接观测证据，因此有人推测它可能是一颗白矮星或其他简并天体$^{[23]}$。该系统距离地球 $550\pm10$ 光年。

Delta Cassiopeiae，也称 Ruchbah 或 Rukbat，意为“膝盖”，可能是一颗 Algol 型食双星，最大亮度约为 2.7 等。有报道称其存在振幅小于 0.1 等、周期约为 $2$ 年 $1$ 个月的食现象$^{[24]}$，但这一结果从未得到确认。它距离地球$99.4\pm0.4$光年$^{[18]}$。Delta Cassiopeiae 曾被让·皮卡尔（Jean Picard）于 1669 年用于测定地球表面上一度弧长，从而推算地球半径$^{[25]}$。

Epsilon Cassiopeiae，又称 Segin，视星等为 3.3，距离地球$410\pm20$光年$^{[18]}$。它是一颗炽热的蓝白色恒星，光谱型为 B3 III，表面温度约为$15{,}680\,\mathrm{K}$。其质量约为太阳的 6.5 倍，直径约为太阳的 4.2 倍，属于 Be 星这一类——高速自转并抛射出环状或壳状物质的恒星$^{[26]}$。

\textbf{较暗恒星}
\begin{figure}[ht]
\centering
\includegraphics[width=6cm]{./figures/f7c9e38a40b34540.png}
\caption{仙后座 $\kappa$ 星及其弓形激波。斯皮策红外望远镜图像（NASA/JPL-Caltech）} \label{fig_Cape_4}
\end{figure}